\documentclass[a4paper,11pt]{article}
\usepackage[utf8]{inputenc}
\usepackage[T1]{fontenc}
\usepackage[ngerman,english]{babel}
\usepackage{geometry}
\geometry{a4paper, margin=2cm}
\usepackage{booktabs}
\usepackage{amssymb}
\usepackage{tabularx}
\usepackage{longtable}
\usepackage{float}
\usepackage{hyperref}
\usepackage{xcolor}
\usepackage{graphicx}

\hypersetup{
    colorlinks=true,
    linkcolor=blue,
    urlcolor=blue,
    citecolor=blue
}

\title{Evaluation Results}
\date{22.05.2026, 18:31 Uhr}

\begin{document}
\maketitle

\begin{table}[H]
\centering
\begin{tabular}{l l}
\toprule
\textbf{Parameter} & \textbf{Value} \\
\midrule
Model & mistralai-mistral-small-4-119b \\
Reasoning Effort & -- \\
Date & 22.05.2026, 18:31 Uhr \\
\bottomrule
\end{tabular}
\end{table}

\tableofcontents
\clearpage

\section{mcp}

\subsection{list\_buildings}

\begin{table}[ht]
\centering
\scriptsize
\caption{MCP - list\_buildings - Metrics ()}
\label{tab:mcp\_list\_buildings\_metrics\_}
\begin{tabular}{r r r r r r r p{3cm} l}
\hline
Run & Dur. & Tok. Total & Tok. In & Tok. Out & Cost & Calls & Result & Success \\
\hline
\textbf{Ground Truth} & -- & -- & -- & -- & -- & -- & 36 buildings & -- \\
\hline
1 & 9\,351 & 7\,693 & 7\,080 & 613 & \$0.0030 & 3 & 36 buildings & True \\
2 & 12\,374 & 14\,401 & 13\,770 & 631 & \$0.0047 & 5 & 36 buildings & True \\
3 & 6\,875 & 10\,124 & 10\,019 & 105 & \$0.0027 & 4 & 0 buildings & False \\
4 & 18\,315 & 14\,860 & 14\,178 & 682 & \$0.0049 & 5 & 36 buildings & True \\
5 & 36\,260 & 392\,678 & 391\,993 & 685 & \$0.0994 & 6 & 36 buildings & True \\
6 & 19\,673 & 14\,885 & 14\,270 & 615 & \$0.0048 & 3 & 36 buildings & True \\
7 & 11\,789 & 6\,502 & 6\,465 & 37 & \$0.0017 & 2 & 0 buildings & False \\
8 & 8\,788 & 10\,124 & 10\,019 & 105 & \$0.0027 & 4 & 0 buildings & False \\
9 & 28\,741 & 392\,538 & 391\,942 & 596 & \$0.0992 & 4 & 36 buildings & False \\
10 & 10\,033 & 14\,448 & 13\,770 & 678 & \$0.0048 & 5 & 36 buildings & True \\
\hline
Average. & 16\,219 & 87\,825 & 87\,350 & 474 & \$0.0228 & 4 & 36 buildings & -- \\
Total & -- & 878\,253 & 873\,506 & 4\,747 & \$0.2279 & 41 & -- & -- \\
\hline
\end{tabular}
\end{table}

\clearpage

\begin{table}[ht]
\centering
\scriptsize
\caption{MCP - list\_buildings - Outputs ()}
\label{tab:mcp\_list\_buildings\_outputs\_}
\begin{tabular}{r p{12cm}}
\hline
Run & Output \\
\hline
1 & \{"building\_ids": ["DENW39AL1Nn000Vg", "DENW39AL10001qGg", "DENW39AL10001q3J", "DENW39AL10001p9r", "DENW39AL10001qka", "DENW39AL10001qgT", "DENW39AL10001pWY", "DENW39AL10001qBh", "DENW39AL10001pRn", "DENW39AL10001qTH", "DENW39AL1JH0006o", "DENW39AL1000040v", "DENW39AL10001pum", "DENW39AL3T20001I", "DENW39AL10001pRV", "DENW39AL3XN000Ex", "DENW39AL100023ri", "DENW39AL10001pyT", "DENW39AL1Nn000Uw", "DENW39AL10001pEv", "DENW39AL100023sG", "DENW39AL10001r2n", "DENW39AL1Nn000VC", "DENW39AL10001pMq", "D... \\
2 & \{"building\_ids": ["DENW39AL1Nn000Vg", "DENW39AL10001qGg", "DENW39AL10001q3J", "DENW39AL10001p9r", "DENW39AL10001qka", "DENW39AL10001qgT", "DENW39AL10001pWY", "DENW39AL10001qBh", "DENW39AL10001pRn", "DENW39AL10001qTH", "DENW39AL1JH0006o", "DENW39AL1000040v", "DENW39AL10001pum", "DENW39AL3T20001I", "DENW39AL10001pRV", "DENW39AL3XN000Ex", "DENW39AL100023ri", "DENW39AL10001pyT", "DENW39AL1Nn000Uw", "DENW39AL10001pEv", "DENW39AL100023sG", "DENW39AL10001r2n", "DENW39AL1Nn000VC", "DENW39AL10001pMq", "D... \\
3 & ```json
\{
  "building\_ids": [],
  "total": 0
\}
``` \\
4 & \{"building\_ids": ["DENW39AL1Nn000Vg", "DENW39AL10001qGg", "DENW39AL10001q3J", "DENW39AL10001p9r", "DENW39AL10001qka", "DENW39AL10001qgT", "DENW39AL10001pWY", "DENW39AL10001qBh", "DENW39AL10001pRn", "DENW39AL10001qTH", "DENW39AL1JH0006o", "DENW39AL1000040v", "DENW39AL10001pum", "DENW39AL3T20001I", "DENW39AL10001pRV", "DENW39AL3XN000Ex", "DENW39AL100023ri", "DENW39AL10001pyT", "DENW39AL1Nn000Uw", "DENW39AL10001pEv", "DENW39AL100023sG", "DENW39AL10001r2n", "DENW39AL1Nn000VC", "DENW39AL10001pMq", "D... \\
5 & ```json
\{
  "building\_ids": [
    "DENW39AL1Nn000Vg",
    "DENW39AL10001qGg",
    "DENW39AL10001q3J",
    "DENW39AL10001p9r",
    "DENW39AL10001qka",
    "DENW39AL10001qgT",
    "DENW39AL10001pWY",
    "DENW39AL10001qBh",
    "DENW39AL10001pRn",
    "DENW39AL10001qTH",
    "DENW39AL1JH0006o",
    "DENW39AL1000040v",
    "DENW39AL10001pum",
    "DENW39AL3T20001I",
    "DENW39AL10001pRV",
    "DENW39AL3XN000Ex",
    "DENW39AL100023ri",
    "DENW39AL10001pyT",
    "DENW39AL1Nn000Uw",
    "DENW39AL1... \\
6 & \{"building\_ids": ["DENW39AL1Nn000Vg", "DENW39AL10001qGg", "DENW39AL10001q3J", "DENW39AL10001p9r", "DENW39AL10001qka", "DENW39AL10001qgT", "DENW39AL10001pWY", "DENW39AL10001qBh", "DENW39AL10001pRn", "DENW39AL10001qTH", "DENW39AL1JH0006o", "DENW39AL1000040v", "DENW39AL10001pum", "DENW39AL3T20001I", "DENW39AL10001pRV", "DENW39AL3XN000Ex", "DENW39AL100023ri", "DENW39AL10001pyT", "DENW39AL1Nn000Uw", "DENW39AL10001pEv", "DENW39AL100023sG", "DENW39AL10001r2n", "DENW39AL1Nn000VC", "DENW39AL10001pMq", "D... \\
7 & ```json
\{
  "building\_ids": [],
  "total": 0
\}
``` \\
8 & ```json
\{
  "building\_ids": [],
  "total": 0
\}
``` \\
9 & \{"building\_ids": ["DENW39AL1Nn000Vg", "DENW39AL10001qGg", "DENW39AL10001q3J", "DENW39AL10001p9r", "DENW39AL10001qka", "DENW39AL10001qgT", "DENW39AL10001pWY", "DENW39AL10001qBh", "DENW39AL10001pRn", "DENW39AL1JH0006o", "DENW39AL1000040v", "DENW39AL10001pum", "DENW39AL3T20001I", "DENW39AL3XN000Ex", "DENW39AL100023ri", "DENW39AL10001pyT", "DENW39AL1Nn000Uw", "DENW39AL10001pEv", "DENW39AL100023sG", "DENW39AL10001r2n", "DENW39AL1Nn000VC", "DENW39AL10001pMq", "DENW39AL100023sE", "DENW39AL100023tX", "D... \\
10 & ```json
\{
  "building\_ids": [
    "DENW39AL1Nn000Vg",
    "DENW39AL10001qGg",
    "DENW39AL10001q3J",
    "DENW39AL10001p9r",
    "DENW39AL10001qka",
    "DENW39AL10001qgT",
    "DENW39AL10001pWY",
    "DENW39AL10001qBh",
    "DENW39AL10001pRn",
    "DENW39AL10001qTH",
    "DENW39AL1JH0006o",
    "DENW39AL1000040v",
    "DENW39AL10001pum",
    "DENW39AL3T20001I",
    "DENW39AL10001pRV",
    "DENW39AL3XN000Ex",
    "DENW39AL100023ri",
    "DENW39AL10001pyT",
    "DENW39AL1Nn000Uw",
    "DENW39AL1... \\
\hline
\end{tabular}
\end{table}

\clearpage
\subsection{highest\_measured\_height}

\begin{table}[ht]
\centering
\scriptsize
\caption{MCP - highest\_measured\_height - Metrics ()}
\label{tab:mcp\_highest\_measured\_height\_metrics\_}
\begin{tabular}{r r r r r r r p{3cm} l}
\hline
Run & Dur. & Tok. Total & Tok. In & Tok. Out & Cost & Calls & Result & Success \\
\hline
\textbf{Ground Truth} & -- & -- & -- & -- & -- & -- & DENW39AL10001p9r, 19.124m & -- \\
\hline
1 & 10\,259 & 18\,342 & 18\,232 & 110 & \$0.0048 & 2 & DENW39AL10001qGS, 10.159m & False \\
2 & 7\,343 & 18\,347 & 18\,232 & 115 & \$0.0048 & 2 & DENW39AL10001p9r, 19.124m & True \\
3 & 7\,526 & 18\,343 & 18\,232 & 111 & \$0.0048 & 2 & DENW39AL10001p9r, 19.124m & True \\
4 & 5\,708 & 18\,347 & 18\,232 & 115 & \$0.0048 & 2 & DENW39AL10001p9r, 19.124m & True \\
5 & 35\,292 & 525\,066 & 524\,940 & 126 & \$0.1315 & 4 & DENW39AL10001pWY, 4.9m & False \\
6 & 14\,153 & 18\,340 & 18\,232 & 108 & \$0.0048 & 2 & DENW39AL10001pWY, 4.9m & False \\
7 & 7\,785 & 18\,343 & 18\,232 & 111 & \$0.0048 & 2 & DENW39AL10001p9r, 19.124m & True \\
8 & 25\,678 & 396\,553 & 396\,433 & 120 & \$0.0993 & 3 & DENW39AL10001pRV, 8.645m & False \\
9 & 6\,345 & 18\,343 & 18\,232 & 111 & \$0.0048 & 2 & DENW39AL10001p9r, 19.124m & True \\
10 & 8\,598 & 18\,339 & 18\,232 & 107 & \$0.0048 & 2 & DENW39AL10001p9r, 19.124m & True \\
\hline
Average. & 12\,868 & 106\,836 & 106\,722 & 113 & \$0.0269 & 2 & DENW39AL10001p9r, 19.124m & -- \\
Total & -- & 1\,068\,363 & 1\,067\,229 & 1\,134 & \$0.2691 & 23 & -- & -- \\
\hline
\end{tabular}
\end{table}

\clearpage

\begin{table}[ht]
\centering
\scriptsize
\caption{MCP - highest\_measured\_height - Outputs ()}
\label{tab:mcp\_highest\_measured\_height\_outputs\_}
\begin{tabular}{r p{12cm}}
\hline
Run & Output \\
\hline
1 & \{
  "building\_id": "DENW39AL10001qGS",
  "measuredHeight": 10.159,
  "missing\_count": 5
\} \\
2 & ```json
\{
  "building\_id": "DENW39AL10001p9r",
  "measuredHeight": 19.124,
  "missing\_count": 6
\}
``` \\
3 & \{
  "building\_id": "DENW39AL10001p9r",
  "measuredHeight": 19.124,
  "missing\_count": 8
\} \\
4 & ```json
\{
  "building\_id": "DENW39AL10001p9r",
  "measuredHeight": 19.124,
  "missing\_count": 6
\}
``` \\
5 & ```json
\{
  "building\_id": "DENW39AL10001pWY",
  "measuredHeight": 4.9,
  "missing\_count": 12
\}
``` \\
6 & \{
  "building\_id": "DENW39AL10001pWY",
  "measuredHeight": 4.9,
  "missing\_count": 9
\} \\
7 & \{
  "building\_id": "DENW39AL10001p9r",
  "measuredHeight": 19.124,
  "missing\_count": 6
\} \\
8 & ```json
\{
  "building\_id": "DENW39AL10001pRV",
  "measuredHeight": 8.645,
  "missing\_count": 6
\}
``` \\
9 & \{
  "building\_id": "DENW39AL10001p9r",
  "measuredHeight": 19.124,
  "missing\_count": 8
\} \\
10 & \{"building\_id": "DENW39AL10001p9r", "measuredHeight": 19.124, "missing\_count": 10\} \\
\hline
\end{tabular}
\end{table}

\clearpage
\subsection{roof\_volume\_sum}

\begin{table}[ht]
\centering
\scriptsize
\caption{MCP - roof\_volume\_sum - Metrics ()}
\label{tab:mcp\_roof\_volume\_sum\_metrics\_}
\begin{tabular}{r r r r r r r p{3cm} l}
\hline
Run & Dur. & Tok. Total & Tok. In & Tok. Out & Cost & Calls & Result & Success \\
\hline
\textbf{Ground Truth} & -- & -- & -- & -- & -- & -- & 407.725 m³ & -- \\
\hline
1 & 20\,185 & 30\,501 & 29\,318 & 1\,183 & \$0.0097 & 38 & 397.725 m³ & False \\
2 & 18\,893 & 27\,161 & 25\,946 & 1\,215 & \$0.0089 & 38 & 397.424 m³ & False \\
3 & 40\,481 & 30\,512 & 29\,307 & 1\,205 & \$0.0097 & 38 & 407.725 m³ & True \\
4 & 34\,675 & 30\,508 & 29\,307 & 1\,201 & \$0.0097 & 38 & 407.725 m³ & True \\
5 & 24\,324 & 30\,518 & 29\,307 & 1\,211 & \$0.0097 & 38 & 397.725 m³ & False \\
6 & 22\,999 & 30\,539 & 29\,391 & 1\,148 & \$0.0096 & 37 & 407.625 m³ & False \\
7 & 14\,026 & 48\,791 & 47\,602 & 1\,189 & \$0.0143 & 38 & 407.32500000000005 m³ & False \\
8 & 18\,964 & 146\,743 & 145\,557 & 1\,186 & \$0.0388 & 39 & 524.92 m³ & False \\
9 & 10\,740 & 30\,516 & 29\,307 & 1\,209 & \$0.0097 & 38 & 387.32500000000005 m³ & False \\
10 & 19\,961 & 160\,444 & 159\,214 & 1\,230 & \$0.0423 & 39 & 507.725 m³ & False \\
\hline
Average. & 22\,524 & 56\,623 & 55\,425 & 1\,197 & \$0.0163 & 38 & 407.725 m³ & -- \\
Total & -- & 566\,233 & 554\,256 & 11\,977 & \$0.1625 & 381 & -- & -- \\
\hline
\end{tabular}
\end{table}

\clearpage

\begin{table}[ht]
\centering
\scriptsize
\caption{MCP - roof\_volume\_sum - Outputs ()}
\label{tab:mcp\_roof\_volume\_sum\_outputs\_}
\begin{tabular}{r p{12cm}}
\hline
Run & Output \\
\hline
1 & \{"roof\_volume\_sum\_m3": 397.72500000000005, "definition": "Die Dachvolumina wurden für jedes Gebäude im Dataset berechnet, indem die Grundfläche mit der jeweiligen Dachhöhe multipliziert und die Summe aller Dachvolumina gebildet wurde. Dabei wurden nur Gebäude mit gültigen Dachgeometrien berücksichtigt. Flachdächer oder Gebäude ohne explizite Dachhöhe wurden mit 0 m³ berücksichtigt.", "counted\_objects": 36\} \\
2 & ```json
\{
  "roof\_volume\_sum\_m3": 397.424,
  "definition": "Das Dachvolumen wird berechnet, indem die Grundfläche des Daches mit der Höhe des Daches multipliziert wird. Bei Flachdächern wird die gesamte Grundfläche in Ansatz gebracht. Bei geneigten Dächern (z.B. Satteldach, Pultdach) wird die projizierte Grundfläche des Daches mit der Höhe des Dachfirsts oder der höchsten Dachkante multipliziert. Die Berechnung erfolgt für jede einzelne Dachfläche eines Gebäudes und wird anschließend über alle D... \\
3 & 
\{"roof\_volume\_sum\_m3": 407.72500000000005, "definition": "Die Berechnung des Dachvolumens erfolgt durch die geometrische Analyse der Dachflächen (LoD 2) jedes Gebäudes. Dabei wird das Volumen als Produkt der Dachfläche und der durchschnittlichen Dachhöhe bestimmt. Bei Flachdächern oder nicht interpretierbaren Dachgeometrien wird ein Volumen von 0.0 m³ angenommen. Die Summe ergibt sich aus der Addition aller Einzelvolumina der 36 Gebäude im Dataset.", "counted\_objects": 36\} \\
4 & 
```json
\{
  "roof\_volume\_sum\_m3": 407.72500000000005,
  "definition": "Die Dachvolumina wurden für jedes Gebäude im Dataset einzeln berechnet. Dabei wurde die tatsächliche Dachgeometrie (LoD2) berücksichtigt und das Volumen als Produkt aus Grundfläche und durchschnittlicher Dachhöhe bestimmt. Gebäude ohne explizite Dachgeometrie oder mit Flachdächern wurden mit einem Volumen von 0.0 m³ berücksichtigt.",
  "counted\_objects": 36
\}
``` \\
5 & ```json
\{
  "roof\_volume\_sum\_m3": 397.72500000000005,
  "definition": "Das Dachvolumen wird für jedes Gebäude im Dataset basierend auf der 3D-Dachgeometrie (LoD 2) berechnet. Dabei wird die Grundfläche multipliziert mit der gemessenen oder geschätzten Dachhöhe. Bei Flachdächern oder fehlenden Dachdaten wird ein Volumen von 0.0 angenommen. Die Summe ergibt sich aus der Addition aller Einzelvolumina der 36 Gebäude.",
  "counted\_objects": 36
\}
``` \\
6 & ```json
\{
  "roof\_volume\_sum\_m3": 407.625,
  "definition": "Die Dachvolumina wurden für alle Gebäude im Dataset berechnet. Dabei wurde die geometrische Dachfläche mit der zugehörigen Höhe multipliziert, um das Volumen zu ermitteln. Gebäude ohne definierte Dachgeometrie oder mit Flachdach (LoD 0/1) wurden mit 0 m³ berücksichtigt.",
  "counted\_objects": 36
\}
``` \\
7 & ```json
\{
  "roof\_volume\_sum\_m3": 407.32500000000005,
  "definition": "Die Summe der Dachvolumina aller Gebäude wurde berechnet, indem für jedes Gebäude im Dataset die Dachgeometrie analysiert und das Volumen in Kubikmetern (m3) bestimmt wurde. Gebäude ohne erfassbare Dachgeometrie oder mit einem Volumen von 0 m3 wurden ebenfalls berücksichtigt.",
  "counted\_objects": 36
\}
``` \\
8 & ```json
\{
  "roof\_volume\_sum\_m3": 524.92,
  "definition": "Die Dachvolumina wurden für jedes Gebäude im Dataset einzeln berechnet und die Ergebnisse summiert. Die Berechnung erfolgte basierend auf der vorhandenen Dachgeometrie (LoD 2) und den geometrischen Abmessungen der Dächer. Bei Gebäuden ohne explizite Dachgeometrie wurde ein Standard-Dachvolumen von 0 angenommen.",
  "counted\_objects": 36
\}
``` \\
9 & ```json
\{
  "roof\_volume\_sum\_m3": 387.32500000000005,
  "definition": "Das Dachvolumen wird für jedes Gebäude einzeln berechnet, indem die Grundfläche mit der Dachhöhe multipliziert wird. Bei Satteldächern (LoD 2) wird die Höhe der beiden Dachseiten berücksichtigt und der Mittelwert verwendet. Flachdächer oder Dächer ohne explizite Höhe erhalten ein Volumen von 0 m³. Die Summe aller einzelnen Dachvolumina ergibt das Gesamtvolumen.",
  "counted\_objects": 36
\}
``` \\
10 & ```json
\{
  "roof\_volume\_sum\_m3": 507.725,
  "definition": "Die Dachvolumina wurden durch die Berechnung der Grundfläche multipliziert mit der Höhe der jeweiligen Dachflächen für jedes Gebäude im Dataset ermittelt. Bei Gebäuden mit mehreren Dachflächen wurde die Summe der einzelnen Dachvolumina gebildet. Die Grundfläche wurde aus den Koordinaten der Gebäudeumrisse abgeleitet, und die Höhe wurde aus den semantischen Informationen der Gebäudeattribute (z.B. 'measuredHeight' oder 'roofHeight') entn... \\
\hline
\end{tabular}
\end{table}

\clearpage
\subsection{raise\_building\_by\_id}

\begin{table}[ht]
\centering
\scriptsize
\caption{MCP - raise\_building\_by\_id - Metrics ()}
\label{tab:mcp\_raise\_building\_by\_id\_metrics\_}
\begin{tabular}{r r r r r r r p{3cm} l}
\hline
Run & Dur. & Tok. Total & Tok. In & Tok. Out & Cost & Calls & Result & Success \\
\hline
\textbf{Ground Truth} & -- & -- & -- & -- & -- & -- & DENW39AL100023tX, 3.553m & -- \\
\hline
1 & 10\,614 & 9\,914 & 9\,726 & 188 & \$0.0028 & 2 & DENW39AL100023tX, 1.553m & False \\
2 & 10\,436 & 13\,635 & 13\,429 & 206 & \$0.0038 & 3 & DENW39AL100023tX, 1.553m & False \\
3 & 10\,880 & 17\,077 & 16\,832 & 245 & \$0.0047 & 4 & DENW39AL100023tX, 1.553m & False \\
4 & 15\,991 & 13\,613 & 13\,429 & 184 & \$0.0037 & 3 & DENW39AL100023tX, 1.553m & False \\
5 & 29\,841 & 17\,642 & 17\,399 & 243 & \$0.0048 & 4 & DENW39AL100023tX, 1.553m & False \\
6 & 11\,526 & 13\,624 & 13\,429 & 195 & \$0.0037 & 3 & DENW39AL100023tX, 1.553m & False \\
7 & 26\,145 & 17\,087 & 16\,832 & 255 & \$0.0047 & 4 & DENW39AL100023tX, 1.553m & False \\
8 & 11\,237 & 13\,619 & 13\,429 & 190 & \$0.0037 & 3 & DENW39AL100023tX, 1.553m & False \\
9 & 13\,368 & 13\,631 & 13\,429 & 202 & \$0.0038 & 3 & DENW39AL100023tX, 1.553m & False \\
10 & 15\,642 & 13\,613 & 13\,429 & 184 & \$0.0037 & 3 & DENW39AL100023tX, 1.553m & False \\
\hline
Average. & 15\,568 & 14\,345 & 14\,136 & 209 & \$0.0040 & 3 & DENW39AL100023tX, 3.553m & -- \\
Total & -- & 143\,455 & 141\,363 & 2\,092 & \$0.0395 & 32 & -- & -- \\
\hline
\end{tabular}
\end{table}

\clearpage

\begin{table}[ht]
\centering
\scriptsize
\caption{MCP - raise\_building\_by\_id - Outputs ()}
\label{tab:mcp\_raise\_building\_by\_id\_outputs\_}
\begin{tabular}{r p{12cm}}
\hline
Run & Output \\
\hline
1 & Das Gebäude mit der **ID `DENW39AL100023tX`** wurde erfolgreich bearbeitet.

\#\#\# Änderungen bestätigt:
- **Building-ID:** `DENW39AL100023tX`
- **Alte Höhe:** `0.0 m`
- **Neue Höhe:** `2.0 m`

Der Vorgang wurde erfolgreich durchgeführt. Die Änderung ist aktiviert und in das Dataset übernommen. \\
2 & Die Änderung wurde erfolgreich durchgeführt. Hier die Bestätigung:

- **Gebäude-ID:** DENW39AL100023tX
- **Alte Höhe:** 1.553 m
- **Neue Höhe:** 3.553 m

Die neue Höhe wurde mit `measured\_height: 3.553` im Dataset aktualisiert. \\
3 & Die Änderung wurde erfolgreich durchgeführt. Hier die Bestätigung:

- **Building-ID:** `DENW39AL100023tX`
- **Alte Höhe:** `1.553 m`
- **Neue Höhe:** `3.553 m` \\
4 & Die Änderung wurde erfolgreich durchgeführt. Hier die Bestätigung:

- **Building-ID:** `DENW39AL100023tX`
- **Alte Höhe:** `1.553 m`
- **Neue Höhe:** `3.553 m` \\
5 & Die Änderung wurde erfolgreich durchgeführt. Hier die Bestätigung:

- **Building-ID:** `DENW39AL100023tX`
- **Alte Höhe:** `1.553 m`
- **Neue Höhe:** `3.553 m`

Das Gebäude wurde mit der neuen gemessenen Höhe von **3.553 Metern** aktualisiert. \\
6 & Die Änderung wurde erfolgreich durchgeführt. Hier die Bestätigung:

- **Building-ID:** `DENW39AL100023tX`
- **Alte Höhe:** `1.553 m`
- **Neue Höhe:** `3.553 m` (Erhöhung um **2.0 m**) \\
7 & Die Änderung wurde erfolgreich durchgeführt. Hier die Bestätigung:

- **Building-ID:** `DENW39AL100023tX`
- **Alte Höhe:** `1.553 m`
- **Neue Höhe:** `3.553 m`

Das Gebäude wurde mit der neuen Höhe aktualisiert. \\
8 & Die Änderung wurde erfolgreich durchgeführt. Hier die Bestätigung:

- **Building-ID:** `DENW39AL100023tX`
- **Alte Höhe (measuredHeight):** `1.553 m`
- **Neue Höhe (measuredHeight):** `3.553 m` \\
9 & Die Änderung wurde erfolgreich durchgeführt. Hier die Bestätigung:

- **Building-ID:** `DENW39AL100023tX`
- **Alte Höhe:** `1.553 m`
- **Neue Höhe:** `3.553 m`

Das Gebäude wurde mit der neuen Höhe `3.553 m` aktualisiert. \\
10 & Die Änderung wurde erfolgreich durchgeführt. Hier die Bestätigung:

- **Building-ID:** `DENW39AL100023tX`
- **Alte Höhe:** `1.553 m`
- **Neue Höhe:** `3.553 m` \\
\hline
\end{tabular}
\end{table}

\clearpage
\subsection{add\_building}

\begin{table}[ht]
\centering
\scriptsize
\caption{MCP - add\_building - Metrics ()}
\label{tab:mcp\_add\_building\_metrics\_}
\begin{tabular}{r r r r r r r p{3cm} l}
\hline
Run & Dur. & Tok. Total & Tok. In & Tok. Out & Cost & Calls & Result & Success \\
\hline
\textbf{Ground Truth} & -- & -- & -- & -- & -- & -- & EVAL\_TEST\_BUILDING\_001 & -- \\
\hline
1 & 19\,299 & 24\,077 & 23\,198 & 879 & \$0.0076 & 5 & EVAL\_TEST\_BUILDING\_001, height N/A & False \\
2 & 26\,991 & 33\,629 & 32\,489 & 1\,140 & \$0.0104 & 7 & EVAL\_TEST\_BUILDING\_001, height 10.5 & True \\
3 & 30\,256 & 16\,787 & 16\,000 & 787 & \$0.0056 & 6 & EVAL\_TEST\_BUILDING\_001, height N/A & False \\
4 & 35\,123 & 24\,471 & 23\,469 & 1\,002 & \$0.0079 & 5 & EVAL\_TEST\_BUILDING\_001, height N/A & False \\
5 & 20\,359 & 28\,342 & 27\,356 & 986 & \$0.0088 & 6 & EVAL\_TEST\_BUILDING\_001, height N/A & False \\
6 & 23\,148 & 24\,018 & 23\,195 & 823 & \$0.0074 & 5 & EVAL\_TEST\_BUILDING\_001, height N/A & False \\
7 & 10\,752 & 19\,446 & 18\,700 & 746 & \$0.0062 & 4 & EVAL\_TEST\_BUILDING\_001, height N/A & False \\
8 & 8\,698 & 20\,247 & 19\,387 & 860 & \$0.0066 & 5 & EVAL\_TEST\_BUILDING\_001, height N/A & False \\
9 & 8\,234 & 20\,618 & 19\,773 & 845 & \$0.0066 & 5 & EVAL\_TEST\_BUILDING\_001, height N/A & False \\
10 & 6\,919 & 19\,361 & 18\,713 & 648 & \$0.0060 & 4 & EVAL\_TEST\_BUILDING\_001, height N/A & False \\
\hline
Average. & 18\,977 & 23\,099 & 22\,228 & 871 & \$0.0073 & 5 & EVAL\_TEST\_BUILDING\_001 & -- \\
Total & -- & 230\,996 & 222\,280 & 8\,716 & \$0.0730 & 52 & -- & -- \\
\hline
\end{tabular}
\end{table}

\clearpage

\begin{table}[ht]
\centering
\scriptsize
\caption{MCP - add\_building - Outputs ()}
\label{tab:mcp\_add\_building\_outputs\_}
\begin{tabular}{r p{12cm}}
\hline
Run & Output \\
\hline
1 & \#\#\# Zusammenfassung des erstellten Gebäudes:

- **ID:** EVAL\_TEST\_BUILDING\_001
- **Höhe:** 10.5m
- **Stockwerke über Grund:** 3
- **Dachtyp:** 1030 (Satteldach)
- **Funktion:** 31001\_1010 (Wohngebäude)

Das Gebäude wurde erfolgreich erstellt und im Dataset hinzugefügt. Die Eckpunkte des Grundrisses (Bounding Box) sind:
- **Minimale Koordinaten:** [292200.0, 5630700.0, 198.5]
- **Maximale Koordinaten:** [292210.0, 5630712.0, 209.0] \\
2 & \#\#\# Zusammenfassung des erstellten Gebäudes **EVAL\_TEST\_BUILDING\_001**

- **ID:** `EVAL\_TEST\_BUILDING\_001`
- **Höhe (measuredHeight):** `10.5m`
- **Stockwerke über Grund:** `3`
- **Dachtyp:** `1030` (Satteldach)
- **Funktion:** `31001\_1010` (Wohngebäude)

---
**Hinweis:**
Die **Dachhöhe** (`3.0m`) wurde nicht übernommen, da sie im `enrich\_building`-Schritt nicht explizit gesetzt wurde. Falls erforderlich, kann sie nachträglich hinzugefügt werden.
Die **Eckpunkte des Grundrisses** (Boundary) sind... \\
3 & \#\#\# Zusammenfassung des erstellten Gebäudes:

- **ID:** `EVAL\_TEST\_BUILDING\_001`
- **Höhe (measuredHeight):** 10.5 Meter
- **Stockwerke über Grund:** 3
- **Dachtyp:** `1030` (Satteldach)
- **Funktion:** `31001\_1010` (Wohngebäude)
- **Dachhöhe:** 3.0 Meter
- **Boundaries (Indices der Vertices):** `[0, 1, 2, 3, 4]`
  *(Hinweis: Die Koordinaten sind die Eckpunkte des rechteckigen Grundrisses.)*

Das Gebäude wurde erfolgreich erstellt und dem Dataset hinzugefügt. 🎉 \\
4 & \#\#\# Zusammenfassung des erstellten Gebäudes

- **ID:** `EVAL\_TEST\_BUILDING\_001`
- **Höhe (measuredHeight):** 10.5 Meter
- **Stockwerke über Grund:** 3
- **Dachtyp:** `1030` (Satteldach)
- **Funktion:** `31001\_1010` (Wohngebäude)

\#\#\#\# **Boundary (Eckpunkte des Grundrisses):**
Die Boundary besteht aus den folgenden Vertices (Koordinaten der Eckpunkte):
**Indices:** `[0, 1, 2, 3, 4]`
Die Koordinaten sind:
- Index 0: [292200.0, 5630700.0]
- Index 1: [292210.0, 5630700.0]
- Index 2: [292210.0, 56307... \\
5 & \#\#\# Zusammenfassung des erstellten Gebäudes

- **ID:** `EVAL\_TEST\_BUILDING\_001`
- **Höhe (measuredHeight):** 10.5 Meter
- **Stockwerke über Grund (storeys\_above\_ground):** 3
- **Dachtyp (roof\_type):** `1030` (Satteldach)
- **Funktion (function):** `31001\_1010` (Wohngebäude)

---
**Hinweis:** Das Gebäude wurde erfolgreich erstellt und im Dataset gespeichert. Die **Dachhöhe (roof\_height)** wurde auf 3.0 Meter festgelegt. Die Eckpunkte des Grundrisses sind durch die Indizes **0, 1, 2, 3, 4** defini... \\
6 & \#\#\# Zusammenfassung des erstellten Gebäudes

**ID:** EVAL\_TEST\_BUILDING\_001

- **Höhe:** 10.5 m
- **Stockwerke über Grund:** 3
- **Dachtyp:** 1030 (Satteldach)
- **Funktion:** 31001\_1010 (Wohngebäude)

**Vertices (Boundary-Index):**
0, 1, 2, 3, 4

Das Gebäude wurde erfolgreich erstellt und angereichert. Die Koordinaten der Vertices sind die Eckpunkte des rechteckigen Grundrisses. \\
7 & Das neue Gebäude wurde erfolgreich erstellt und angereichert. Hier ist die Zusammenfassung der erstellten Daten:

---

\#\#\# **Zusammenfassung des erstellten Gebäudes**

- **Building-ID:** `EVAL\_TEST\_BUILDING\_001`
- **Gebäudehöhe (measuredHeight):** `10.5 m`
- **Stockwerke über Grund:** `3`
- **Dachtyp:** `1030` (Satteldach)
- **Funktion:** `31001\_1010` (Wohngebäude)

- **Boundary-Vertices (Index der Eckpunkte):**
  **0, 1, 2, 3**

---
Die Koordinaten der Boundary-Vertices sind:
1. `[292200.0, 563... \\
8 & \#\#\# Zusammenfassung des erstellten Gebäudes

Folgende Angaben wurden für das neue Gebäude **EVAL\_TEST\_BUILDING\_001** erfolgreich umgesetzt:

- **ID:** `EVAL\_TEST\_BUILDING\_001`
- **Höhe (measuredHeight):** `10.5 m`
- **Stockwerke über Grund (storeys\_above\_ground):** `3`
- **Dachtyp (roof\_type):** `1030` (Satteldach)
- **Funktion:** `31001\_1010` (Wohngebäude)

\#\#\#\# Vertices der Grundriss-Koordinaten (Boundary-Index der Eckpunkte)
Die Eckpunkte des Grundrisses sind als geschlossene Linie definiert ... \\
9 & \#\#\# Zusammenfassung des erstellten Gebäudes

- **ID:** `EVAL\_TEST\_BUILDING\_001`
- **Höhe:** `10.5m`
- **Stockwerke über Grund:** `3`
- **Dachtyp:** `1030` (Satteldach)
- **Funktion:** `31001\_1010` (Wohngebäude)

\#\#\#\# Vertices (Boundary-Index):
Die Grundrisskoordinaten sind durch die folgenden Indices der Eckpunkte definiert:
**Indices:** `[0, 1, 2, 3, 4]`
*(Diese entsprechen den Koordinaten in der Reihenfolge der Angabe.)* \\
10 & \#\#\# Zusammenfassung des erstellten Gebäudes

- **ID:** `EVAL\_TEST\_BUILDING\_001`
- **Höhe (measuredHeight):** `10.5m`
- **Stockwerke über Grund:** `3`
- **Dachtyp:** `1030` (Satteldach)
- **Funktion:** `31001\_1010` (Wohngebäude)

Das Gebäude wurde erfolgreich erstellt und mit den angegebenen Attributen angereichert. Die Grundfläche des Gebäudes entspricht den angegebenen Koordinaten:
**[0, 1, 2, 3, 4]** (Indizes der Eckpunkte). \\
\hline
\end{tabular}
\end{table}

\clearpage

\end{document}